\documentclass[12pt,a4paper]{article}
\usepackage[utf8]{inputenc}
\usepackage[T1]{fontenc}
\usepackage{amsmath,amssymb}
\usepackage{enumitem}
\usepackage{geometry}
\usepackage{array}
\usepackage{multicol}
\geometry{margin=2.2cm}
\setlist{itemsep=0.25em, topsep=0.3em}
\title{Aula 11 -- Técnicas de Programação SQL}
\author{Banco de Dados -- Guia de Revisão}
\date{}
\begin{document}
\maketitle
\tableofcontents
\newpage

\section{Objetivo da aula}
Esta aula mostra como programas acessam bancos de dados. O foco é entender a diferença entre SQL interativo e SQL usado dentro de aplicações, além de JDBC, SQL embutido, cursores, Hibernate, procedures e SQL Injection.

\section{Ideia central}
Na prática, usuários raramente digitam SQL diretamente. Aplicações executam SQL por meio de uma linguagem hospedeira, como Java, C, C++ ou Python.

\section{Linguagem hospedeira e sublinguagem de dados}
\begin{itemize}
    \item \textbf{Linguagem hospedeira}: linguagem geral da aplicação, como Java.
    \item \textbf{Sublinguagem de dados}: SQL usado para acessar o banco.
\end{itemize}

\section{Descompasso de impedância}
Ocorre porque linguagens de programação e bancos relacionais representam dados de formas diferentes.

Exemplo:
\begin{itemize}
    \item Java usa objetos, listas, tipos, referências.
    \item SQL retorna tabelas, linhas e colunas.
\end{itemize}

Por isso são necessários cursores, iteradores, \texttt{ResultSet} ou frameworks ORM.

\section{SQL embutido}
SQL embutido coloca comandos SQL dentro de uma linguagem hospedeira. Em C, por exemplo, comandos podem aparecer com \texttt{EXEC SQL}.

Conceitos importantes:
\begin{itemize}
    \item Variáveis compartilhadas entre programa e SQL.
    \item Pré-processador para separar comandos SQL.
    \item \texttt{SQLCODE} ou \texttt{SQLSTATE} para erros.
    \item Cursor para resultados com várias linhas.
\end{itemize}

\section{Cursores}
Cursores permitem processar resultado linha por linha.

Fluxo típico:
\begin{verbatim}
DECLARE cursor
OPEN cursor
FETCH cursor
CLOSE cursor
\end{verbatim}

\texttt{FETCH} avança para a próxima tupla.

\section{SQL dinâmica}
SQL dinâmica permite montar comandos em tempo de execução. É útil quando a consulta depende de escolhas do usuário, mas aumenta riscos se for feita com concatenação insegura.

\section{JDBC}
JDBC é uma API Java para acesso a banco de dados.

Fluxo comum:
\begin{enumerate}
    \item Abrir \texttt{Connection}.
    \item Criar \texttt{Statement} ou \texttt{PreparedStatement}.
    \item Executar comando.
    \item Processar \texttt{ResultSet}, se houver.
    \item Fechar recursos.
\end{enumerate}

\subsection{Exemplo visual: consulta segura}
\begin{verbatim}
Connection con = DriverManager.getConnection(url, user, pass);
PreparedStatement ps = con.prepareStatement(
    "SELECT * FROM Aluno WHERE matricula = ?"
);
ps.setInt(1, matricula);
ResultSet rs = ps.executeQuery();
while (rs.next()) {
    String nome = rs.getString("nome");
}
\end{verbatim}

\begin{itemize}
    \item \texttt{Connection}: conexão com o banco.
    \item \texttt{PreparedStatement}: comando parametrizado.
    \item \texttt{?}: parâmetro.
    \item \texttt{setInt(1,...)}: preenche o primeiro parâmetro.
    \item \texttt{ResultSet}: linhas retornadas.
\end{itemize}

\section{executeQuery e executeUpdate}
\begin{itemize}
    \item \texttt{executeQuery()}: usado para \texttt{SELECT}; retorna \texttt{ResultSet}.
    \item \texttt{executeUpdate()}: usado para \texttt{INSERT}, \texttt{UPDATE}, \texttt{DELETE}; retorna quantidade de linhas afetadas.
\end{itemize}

\section{Transações em JDBC}
Para controlar transação manualmente:
\begin{verbatim}
con.setAutoCommit(false);
try {
    // comandos SQL
    con.commit();
} catch (Exception e) {
    con.rollback();
}
\end{verbatim}

\section{SQL Injection}
SQL Injection ocorre quando entrada externa é concatenada diretamente no SQL.

Código vulnerável:
\begin{verbatim}
String sql = "SELECT * FROM Usuario WHERE login='" + login +
             "' AND senha='" + senha + "'";
\end{verbatim}

Entrada perigosa:
\begin{verbatim}
' OR '1'='1
\end{verbatim}

Correção: usar \texttt{PreparedStatement} com parâmetros.

\section{Hibernate e ORM}
Hibernate é um framework de mapeamento objeto-relacional. Ele associa classes Java a tabelas relacionais.

Exemplo conceitual:
\begin{center}
\begin{tabular}{|c|c|}
\hline
Java & Banco Relacional \\
\hline
Classe \texttt{Aluno} & Tabela \texttt{Aluno} \\
Objeto & Linha \\
Atributo do objeto & Coluna \\
\hline
\end{tabular}
\end{center}

\section{Procedures, funções e SQL/PSM}
Procedures e funções armazenadas ficam no SGBD e podem ser chamadas por aplicações.

Vantagens:
\begin{itemize}
    \item centralizam lógica perto dos dados;
    \item reduzem repetição de código;
    \item podem melhorar controle e padronização.
\end{itemize}

\section{Senhas e dados sensíveis}
Senhas não devem ser armazenadas em texto puro. Prática melhor: hash com salt.

Fluxo:
\begin{enumerate}
    \item Usuário digita senha.
    \item Sistema combina senha com salt.
    \item Calcula hash.
    \item Compara com hash armazenado.
\end{enumerate}

\section{Pegadinhas comuns}
\begin{itemize}
    \item \texttt{PreparedStatement} não serve só para performance; também ajuda contra SQL Injection.
    \item \texttt{ResultSet} começa antes da primeira linha; precisa de \texttt{rs.next()}.
    \item \texttt{executeQuery} é para consulta; \texttt{executeUpdate} é para alteração.
    \item Hibernate não elimina o modelo relacional; ele mapeia objetos para relações.
\end{itemize}

\section{Desafios}
\subsection*{Desafio 1}
Explique por que o código abaixo é vulnerável:
\begin{verbatim}
String sql = "DELETE FROM Employee WHERE ssn=" + ssn;
\end{verbatim}

\subsection*{Desafio 2}
Qual método JDBC você usaria para executar um \texttt{SELECT}? E para executar um \texttt{UPDATE}?

\section{Resolução dos desafios}
\textbf{Desafio 1}: concatena entrada diretamente no SQL. Um valor malicioso pode alterar a condição e afetar mais linhas que o esperado.

\textbf{Desafio 2}: \texttt{executeQuery()} para \texttt{SELECT}; \texttt{executeUpdate()} para \texttt{UPDATE}, \texttt{INSERT} e \texttt{DELETE}.

\section{Checklist final}
\begin{itemize}
    \item Sei explicar descompasso de impedância?
    \item Sei a função de \texttt{Connection}, \texttt{PreparedStatement} e \texttt{ResultSet}?
    \item Sei reconhecer SQL Injection?
    \item Sei diferenciar cursor, JDBC, SQLJ, Hibernate e stored procedure?
\end{itemize}

\end{document}
